\exercisesfor{1}

\begin{enumerate}
\def\labelenumi{\arabic{enumi}.}
\item
  假设 K 是特征为 0 的域。令 E 为 K 上有限秩的余交换双代数。证明 \(P(E) = \{0\}\)（应用第 6 号定理 1）。
\item
  令 E 为满足 \(\mathrm{P}(\mathrm{E})=\{0\}\) 的双代数，并令 \((\mathrm{E}_{n})_{n\geqslant0}\) 为与其双代数结构相容的 E 的一个滤子。证明对 n 作归纳可得，对所有 \(E_{n}^{+}=\{0\}\) 有 \(n\geqslant0\)，并由此推出 \(E^{+}=\{0\}\)，即 E 退化为 K。
\item
  令 \(G\) 为幺半群，\(E = K[G]\) 为其代数（《代数》，第 III 章，§2，第 6 号）。
\end{enumerate}

\begin{enumerate}
\def\labelenumi{(\alph{enumi})}
\item
  证明 E 上存在唯一的余代数结构，使得对所有 \(c(g) = g \otimes g\) 有 \(g \in G\)；这个结构与 E 上的代数结构相容，并使 E 成为余交换双代数，其余单位 \(\varepsilon\) 满足对所有 \(\varepsilon(g) = 1\) 有 \(g \in G\)。
\item
  证明 E 的每个本原元素都为零。由此推出，若 \(G \neq \{e\}\)，则双代数 E 不存在与其双代数结构相容的滤子（应用习题 2）。
\item
  假设 \(\mathbf{K}\) 是整环。证明，\(\mathrm{G}\) 中的元素是满足 \(x\in \mathrm{E}\) 且 \(c(x) = x\otimes x\) 的唯一非零元素。证明，\(\mathrm{E}\) 有逆运算 \(i\)（参见《代数》，第 III 章，§ 11，习题 4）当且仅当 \(\mathrm{G}\) 是群，并且在该情形下 \(i(g) = g^{-1}\) 对所有 \(g\in \mathrm{G}\) 成立。
\end{enumerate}

\begin{enumerate}
\def\labelenumi{\arabic{enumi}.}
\setcounter{enumi}{3}
\item
  令 \(\mathbf{E}\) 为余交换双代数，令 \(\mathbf{G}\) 为满足 \(g \in \mathbf{E}\) 且 \(\varepsilon(g) = 1\) 的 \(c(g) = g \otimes g\) 所成的集合。证明 \(\mathbf{G}\) 对乘法封闭，并证明当 \(\mathbf{E}\) 有逆运算时它是一个群。证明，若 \(\mathbf{K}\) 是域，则 \(\mathbf{G}\) 的元素在 \(\mathbf{K}\) 上线性无关。
\item
  令 E 为双代数，令 \(E' = \text{Hom}(E, K)\) 为其对偶，并赋予它由 E 上余代数结构通过对偶性导出的代数结构（《代数》，第 III 章，§11，第 1 号）。令 m 为从 \(u \mapsto u(1)\) 到 \(\mathbf{E}'\) 的同态 \(\mathbf{K}\) 的核；它是 \(\mathbf{E}'\) 的理想。证明，若 \(u \in \mathfrak{m}^2\) 且 \(x \in \mathrm{P}(\mathrm{E})\)，则 \(u(x) = 0\)。当 \(\mathbf{K}\) 是域时，证明 \(\mathrm{P}(\mathrm{E})\) 是 \(\mathfrak{m}^2\) 中 \(\mathbf{E}^+\) 的正交补；推出 \(\dim \mathrm{P}(\mathrm{E}) \leqslant \dim \mathfrak{m} / \mathfrak{m}^2\)，并证明当 \(\mathbf{E}\) 在 \(\mathbf{K}\) 上有限秩时取等号。
\item
  令 E 为双代数，令 \(u \in \mathrm{E}' = \operatorname{Hom}(\mathrm{E}, \mathrm{K})\)（参见习题 5）。令 \(g \in \mathrm{E}\) 满足 \(\varepsilon(g) = 1\) 且 \(c(g) = g \otimes g\)。证明映射 \(u \mapsto u(g)\) 是从 \(\mathrm{E}'\) 到 \(\mathrm{K}\) 的代数同态，且映射 \(y \mapsto gy\) 是余代数 \(\mathrm{E}\) 的自同态；证明当 \(\mathrm{E}\) 有逆运算时，这个自同态是自同构。
\item
  假设 K 是域。令 E 为满足以下条件的双代数：
\end{enumerate}

\begin{enumerate}
\def\labelenumi{(\roman{enumi})}
\item
  E 是余交换的；
\item
  E 有逆运算（《代数》，第 III 章，§ 11，习题 4）；
\item
  \(\mathrm{P(E)} = \{0\}\)；
\item
  E 在 K 上有限秩。
\end{enumerate}

令 \(\mathbf{E}'\) 表示余代数 E 的对偶 K-代数。令 G 表示满足 \(g \in \mathbf{E}\) 且 \(\varepsilon(g) = 1\) 的 \(c(g) = g \otimes g\) 所成的集合；它是一个群（参见习题 4）。

\begin{enumerate}
\def\labelenumi{(\alph{enumi})}
\item
  令 \(g \in \mathbf{G}\)，令 \(\mathfrak{m}_g\) 为从 \(u \mapsto u(g)\) 到 \(\mathbf{E}'\) 的同态 \(\mathbf{K}\) 的核。证明 \(\mathfrak{m}_g = \mathfrak{m}_g^2\)（使用习题 6 将其归结为 \(g = 1\) 的情形；然后应用习题 5）。
\item
  假设 K 代数闭。证明理想 \(m_{g}\) 是 \(E'\) 的唯一极大理想；推出它们的交为 \(\{0\}\)，\(E'\) 可认作乘积 \(K^{G}\)，并且 E 可认作有限群 G 的双代数 K[G]（参见习题 3）。
\end{enumerate}

\begin{enumerate}
\def\labelenumi{\arabic{enumi}.}
\setcounter{enumi}{7}
\item
  证明李代数 \(\mathfrak{g}\) 的包络双代数有逆运算 \(i\)，且 \(i(x) = -x\) 对所有 \(x \in \sigma(\mathfrak{g})\) 成立。
\item
  假设 K 是 Q-代数。令 g 为具有 K-基的李代数，令 U 为其带有典范滤子 \((\mathrm{U}_{n})_{n\geqslant0}\) 的包络双代数。令 \(x\in U^{+}\)，令 n 为一个 \(\geqslant1\) 的整数。证明，x 属于 \(U_{n}^{+}\) 当且仅当 \(c^{+}(x)\) 属于 \(\sum_{\substack{i+j=n\\ i,j\geqslant1}}\mathrm{Im}(\mathrm{U}_{i}^{+}\otimes\mathrm{U}_{j}^{+})\)。
\item
  假设 K 是 Q-代数。令 E 为带有与其双代数结构相容的滤子的余交换 K-双代数。证明态射
\end{enumerate}

\[
f _ {\mathbf {E}} \colon \mathrm{U} (\mathrm{P} (\mathrm{E})) \rightarrow \mathrm{E}
\]

在第 4 号命题 8 中定义的态射是同构，只要假设 P(E) 是自由 K-模（证明与定理 1 相同）。

¶ 11. 设 E 为双代数，I 为有限集。给定一组基 \((e_{\alpha})\)

它由 \(\alpha\) 中的元素 \(\mathbf{N}^{\mathrm{I}}\) 编号，并满足：

\begin{enumerate}
\def\labelenumi{(\roman{enumi})}
\item
  \(e_{0} = 1\)；
\item
  对所有 \(c(e_{\gamma}) = \sum_{\alpha + \beta = \gamma} e_{\alpha} \otimes e_{\beta}\)，有 \(\gamma \in \mathbf{N}^{\mathrm{I}}\)。
\end{enumerate}

条件 (ii) 蕴含，作为余代数，E 同构于 \(\mathrm{TS}(\mathrm{K}^{\mathrm{I}})\) 的对称张量余代数 \(\mathrm{K}^{\mathrm{I}}\)（参见《代数》，第 IV 章，§ 5，第 7 号）。

\begin{enumerate}
\def\labelenumi{(\alph{enumi})}
\tightlist
\item
  令 \(\mathrm{E}' = \mathrm{Hom}(\mathrm{E},\mathrm{K})\) 为 \(\mathbf{E}\) 的对偶代数，令 \(\Lambda = \mathrm{K}[[(\mathrm{X}_i)_{\mathrm{i}\in I}]]\) 为关于由 I 编号的不定元 \(\mathbf{X}_i\) 的形式幂级数代数。对所有 \(u\in \mathrm{E}'\)，令 \(\phi_u\) 为形式幂级数
\end{enumerate}

\[
\sum_ {\alpha} u (e _ {\alpha}) x ^ {\alpha},
\]

其中 \(x^{\alpha} = \prod_{i\in I}\mathrm{X}_{i}^{\alpha (i)}\)

证明 \(u \mapsto \phi_u\) 是从 \(E'\) 到 \(\Lambda\) 的同构，并且这个同构将 \(E'\) 上的逐点收敛拓扑映到 \(\Lambda\) 上的乘积拓扑。

\begin{enumerate}
\def\labelenumi{(\alph{enumi})}
\setcounter{enumi}{1}
\tightlist
\item
  令 \(c_{\alpha \beta \gamma}\) 为代数 \(E\) 关于基 \((e_{\alpha})\) 的结构常数。则
\end{enumerate}

\[
e _ {\alpha} e _ {\beta} = \sum_ {\gamma} c _ {\alpha \beta \gamma} e _ {\gamma}.
\]

用 x 代替 \((\mathbf{X}_{i})_{i\in\mathbb{I}}\)，用 y 代替 \((\mathbf{Y}_{i})_{i\in\mathbb{I}}\)，其中 \(Y_{i}\) 是新的不定元。证明，存在关于变量 x、y 的形式幂级数族 \(f(\boldsymbol{x},\boldsymbol{y})=(f_{i}(\boldsymbol{x},\boldsymbol{y}))_{i\in\mathbb{I}}\)，使得

\[
\boldsymbol {f} (\boldsymbol {x}, \boldsymbol {y}) ^ {\gamma} = \sum_ {\alpha , \beta} c _ {\alpha \beta \gamma} \boldsymbol {x} ^ {\alpha} \boldsymbol {y} ^ {\beta} \quad \text { for   all } \gamma \in \mathbf {N} ^ {\mathrm{I}},
\]

其中如上记 \(x^{\alpha} = \prod_{i} X_{i}^{\alpha(i)}\)，对于 \(y^{\beta}\) 和 \(f(x, y)^{\gamma}\) 也作类似记号。

证明 \(f(x, y)\) 是 \(K\) 上维数为 \(n = \text{Card}(I)\) 的形式群律，其意义见第 I 章 § 1 习题 24†；定义这个形式群律到李代数 \(P(E)\) 的李代数同构，其中 \(e_{\alpha}\)（满足 \(|\alpha| = 1\)）构成一组基。

\begin{enumerate}
\def\labelenumi{(\alph{enumi})}
\setcounter{enumi}{2}
\item
  反之，证明 K 上关于 x、y 的每个形式群律都可由上述步骤得到，并且所得结果在同构意义下唯一。
\item
  当 K 是 Q-代数时，将上述结果应用于在 K 上具有有限基的李代数 g 的包络双代数。推出存在以 g 为李代数的形式群律（且在同构意义下唯一）。
\item
  明确给出与单参数形式群律 \(f(x, y) = x + y\)（分别为 \(f(x, y) = x + y + xy\)）相对应的双代数。
\end{enumerate}

¶ 12. 假设 K 是特征为 p > 0 的域。

\begin{enumerate}
\def\labelenumi{(\alph{enumi})}
\tightlist
\item
  令 \(\mathfrak{g}\) 为李 \(p\)-代数（第 I 章，§ 1，习题 20），令 \(\tilde{\mathbf{U}}\) 为其
\end{enumerate}


限制包络双代数（第 I 章，§2，习题 6）。证明 \(\tilde{\mathbf{U}}\) 上存在唯一的双代数结构，它与其代数结构相容，并且 \(g\) 的元素是本原元。证明 \(\tilde{\mathbf{U}}\) 是余交换的，且 \(\mathrm{P}(\tilde{\mathbf{U}}) = g\)。

\begin{enumerate}
\def\labelenumi{(\alph{enumi})}
\setcounter{enumi}{1}
\item
  若 \(n\) 是一个 \(\geqslant 0\) 的整数，令 \(\tilde{\mathbf{U}}_n\) 表示 \(\tilde{\mathbf{U}}\) 中由乘积 \(x_1\cdots x_n\) 生成的向量子空间，其中 \(x_i \in \mathfrak{g}\) 对所有 \(i\) 成立。证明 \((\tilde{\mathbf{U}}_n)_{n \geqslant 0}\) 是与 \(\tilde{\mathbf{U}}\) 的双代数结构相容的滤子。
\item
  令 \(\mathbf{E}\) 为双代数，令 \(\mathfrak{g} = \mathrm{P}(\mathbf{E})\) 为其本原元素的李代数。映射 \(x\mapsto x^p\) 保持 \(\mathfrak{g}\) 不变，并赋予 \(\mathfrak{g}\) 李 \(p\)-代数结构。证明典范嵌入 \(\mathfrak{g}\to \mathbf{E}\) 可唯一延拓为双代数态射 \(\tilde{\mathbf{U}}\rightarrow \mathbf{E}\)，且该态射是单射（利用以下事实：若 \((x_{i})_{i\in I}\) 是带有全序的 \(\mathfrak{g}\) 的一组基，则单项式 \(\prod_{i\in I}x_i^{n_i}(0\leqslant n_i < p)\) 构成 \(\mathbf{U}\) 的一组基（同上），并仿照引理 2 的证明进行论证）。
\end{enumerate}
% semantic-fragments: legacy-input-seam
\relax{}
